% LTeX: language=en-GB
% !TEX root = ../../report.tex
\section{Step 4 - Design variables optimisation}\label{section:optimisation-code}
% objective of the step; inputs inherited:
%   impact velocity (Step 0, chosen percentile), acceleration allowable (Step 3, x2 impact)
% output: optimal design variables + their sensitivity -> design recommendations

\subsection{Implementation}

\subsubsection{Problem formulation}
% design variables + bounds (e.g. airbag dimensions/volume, vent area, inflation pressure)
% objective (e.g. airbag mass)
% constraints: peak payload acceleration <= allowable; no bottoming-out;
%   fabric/seam stress <= seam efficiency x fabric strength
% fixed parameters and their source (Step 0, Step 3, datasheets)
* codice ispirato da Zhou et al 2019 - spiegazione del modello e dell'aspetto dei vari blocchi
* Sapendo che le thresholds delle provette sono tot, allora si può convergere verso un pre-dimensionamento dell'airbag con tot dimensioni o caratteristiche. 
* oracle code for direzionamento verso un punto di partenza per il design

* I CHOSE THE CRUDE IMPLEMENTATION --- THIS IS BECAUSE THE ONLY COUPLING ANALYTICAL MODEL IS MISSING AND I CANNOT RETRIEVE IT. THEREFORE THIS IS MY IMPLEMENTATION XXX
Title: "Design, Fabrication and Testing of an Airdrop Platform Utilizing Airbags as Shock Absorbers"

Author: Timothy F. Patterson

Report: NATICK/TR-87/042 (DTIC accession ADA199593)

Date: August 1985 (work performed March 1984–May 1985)


\subsubsection{Physical model}
% 1DoF impact model -> extension to 2D (horizontal velocity component)
% gas model (thermodynamic assumption), venting
% fabric transpiration (porosity leakage)
% seam efficiency in the membrane stress check
\begin{lstlisting}[language=Python, frame=single,
  literate={·}{{$\cdot$}}1 {→}{{$\rightarrow$}}2 {−}{{-}}1
           {ρ}{{$\rho$}}1 {φ}{{$\varphi$}}1 {ṁ}{{$\dot{m}$}}1
           {≤}{{$\leq$}}1 {≥}{{$\geq$}}1]
init: geometry(h=0), m0 = P0·V0/(R·T0)
loop while t < t_max:
  choose dt from previous load (hysteresis)
  F_n = max((P−Pamb)·A_prev, 0) → a_x, a_y          [explicit: old P, old A]
  u_new = u + a·dt ;  h_new = clamp(h + u_n,new·dt)  [semi-implicit Euler]
  A, V = geometry(h_new) ;  A_vent = A_or + φ(S − A)
  ρ = m/V → P, T adiabatic from previous state
  repeat N sub-steps (dt/N):                          [sub-cycling, not iteration]
      if P ≤ Pamb: stop
      dm = ṁ(P,T,A_vent)·dt/N → m, ρ, P (old T), T
  P = max(P, Pamb)
  if P > P_burst → BURST  | if h ≥ D0 → BOTTOMED | if y_pen ≤ 0 → REBOUND
  advance state
\end{lstlisting}


\subsubsection{Optimisation algorithm}
% differential evolution: why gradient-free (discontinuous constraints, bottoming-out)
% settings: population, generations, tolerance, seed
% constraint handling (penalty / feasibility rule)
\begin{table}[h!]
\centering
\caption{Inputs of the preliminary airbag optimisation: fixed inputs and boundary conditions (left) and optimisable design variables (right).}
\label{tab:airbag-design-variables-selected}
\begin{thesistable}{
  colspec = {Q[l] X[1,l] Q[l] X[1.12,l]},
  % ---- column headers (blue = fixed inputs, teal = design variables)
  cell{1}{1} = {c=2}{c, bg=phaseA},
  cell{1}{3} = {c=2}{c, bg=phaseB},
  vline{3} = {wd=\gap, fg=white},
  % ---- category bands, left half
  cell{2,7,11}{1} = {c=2}{bg=tintA, fg=phaseA, font=\tblfont\sffamily\bfseries},
  hline{7,11} = {1-2}{0.6pt, fg=grouprulecol},
  % ---- category bands, right half
  cell{2,7,9}{3} = {c=2}{bg=tintB, fg=phaseB, font=\tblfont\sffamily\bfseries},
  hline{7,9} = {3-4}{0.6pt, fg=phaseB!45},
  % ---- derived quantities (not inputs), greyed out
  cell{6,14}{1-2} = {fg=navy!60, font=\tblfont\itshape},
  % ---- categorical variables: one row per option
  cell{11}{3} = {r=3}{l},
  cell{14}{3} = {r=5}{l},
}
Fixed inputs and boundary conditions & & Optimisable design variables & \\
% ------------------------------------------------------------------------------------------------
Impact scenario & & Continuous (searched by the optimiser) & \\
$M_\text{p}$  & \qty{0.5}{\kilogram}, payload mass
  & $D_0$ & \qtyrange{0.39}{1.5}{\metre}, diameter \\
$u_0$  & \qty{17.2}{\metre\per\second}, vertical impact speed
  & $L_0/D_0$ & \numrange{0.5}{4.0}, aspect ratio (cylinder only) \\
$u_{x,0}$ & \qty{5.0}{\metre\per\second}, horizontal impact speed
  & $P_0$ & $P_\text{amb}$ to \qty{1.4e5}{\pascal}, inflation pressure \\
$u_{n,0}$ & $\approx$\,\qty{17.9}{\metre\per\second} along $\mathbf{n}$, at $\theta \approx \ang{16.2}$ (derived)
  & $\varphi_\text{vent}$ & \numrange{0}{0.10}, vent open-area fraction \\
% ------------------------------------------------------------------------------------------------
Environment & & Discrete -- integer (searched by the optimiser) & \\
$P_\text{amb}$ & \qty{101325}{\pascal}, ambient pressure
  & $n_\text{ply}$ & $\{1,\dots,4\}$, number of fabric plies \\
$g$ & \qty{9.81}{\metre\per\second\squared}
  & Discrete -- categorical (enumerated) & \\
$T_0$ & \qty{288}{\kelvin}, initial gas temperature
  & Shape & sphere ($k = 0.5$), cylinder ($k = 1$) \\
% ------------------------------------------------------------------------------------------------
Requirements and safety factors & & Fabric & Nylon 6,6, silicone-coated: \qty{600}{\newton\per\centi\metre}, \qty{180}{\gram\per\metre\squared} \\
$a_\text{allow}$ & \qty{500}{\metre\per\second\squared}, allowable payload acceleration
  & & PET, PU-coated: \qty{550}{\newton\per\centi\metre}, \qty{190}{\gram\per\metre\squared} \\
$\mathit{SF}$ & 1.2, safety factor on $a_\text{allow}$
  & & Aramid, coated: \qty{1000}{\newton\per\centi\metre}, \qty{260}{\gram\per\metre\squared} \\
$a_\text{lim}$ & $a_\text{allow}/\mathit{SF} \approx \qty{417}{\metre\per\second\squared}$ (derived)
  & Gas & N\textsubscript{2}: $R = \qty{296.8}{\joule\per\kilogram\per\kelvin}$, $\gamma = 1.400$ \\
$u_\text{res,max}$ & \qty{2.0}{\metre\per\second}, residual speed cap
  & & He: $R = \qty{2076.9}{\joule\per\kilogram\per\kelvin}$, $\gamma = 1.667$ \\
$\eta_\text{seam}$ & 0.70, seam efficiency on fabric strength
  & & Ar: $R = \qty{208.1}{\joule\per\kilogram\per\kelvin}$, $\gamma = 1.667$ \\
$\mathit{SF}_\text{burst}$ & 1.50, design factor on membrane tension
  & & CO\textsubscript{2}: $R = \qty{188.9}{\joule\per\kilogram\per\kelvin}$, $\gamma = 1.289$ \\
$\mathit{PEN}$ & $10^4$, penalty weight on constraint violations
  & & Dry air: $R = \qty{287.0}{\joule\per\kilogram\per\kelvin}$, $\gamma = 1.400$ \\
\end{thesistable}
\end{table}



\subsection{Results}
\subsubsection{Reference optimum}         % design variables, objective, which constraints are active
\subsubsection{Convergence of the optimiser} % objective vs generation, repeat with different seeds
\subsubsection{Effect of the 2D impact}  % vertical-only vs with horizontal velocity
\subsubsection{Effect of transpiration and seam efficiency}
\subsubsection{Sensitivity to the inherited inputs} % allowable, impact velocity, payload mass
% -> closing paragraph bridges to the design recommendations chapter (payload-specific)

\subsubsection{Step 4 deliverable: airbag sizing parameters}





